% Supplementary material for main.tex (MethodsX). Build with build.sh: tables and sections are referenced from main.tex through xr.
\documentclass[11pt]{article}
\usepackage[a4paper,margin=25mm]{geometry}
\usepackage[T1]{fontenc}
\usepackage{lmodern}
\usepackage{microtype}
\usepackage[british]{babel}
\usepackage{amsmath,amssymb}
\usepackage{array,booktabs,tabularx,multirow,makecell}
\usepackage{graphicx}
\usepackage{xcolor}
\usepackage[numbers,sort&compress]{natbib}
\usepackage{xr-hyper}
\usepackage[hidelinks]{hyperref}
% Shared by main.tex and supplement.tex.
\newcommand{\todo}[1]{\textcolor{red}{[TODO: #1]}}
\newcommand{\RePV}{RePV}            % method name
% Coverage symbols of the tool comparison table (main.tex): full, partial, absent. Needs tikz where typeset.
\DeclareRobustCommand{\cF}{\tikz[baseline=-0.6ex]\fill (0,0) circle (0.55ex);}
\DeclareRobustCommand{\cH}{\tikz[baseline=-0.6ex]{\draw[line width=0.4pt] (0,0) circle (0.55ex); \fill (0,0) -- (90:0.55ex) arc (90:270:0.55ex) -- cycle;}}
\DeclareRobustCommand{\cN}{\tikz[baseline=-0.6ex]\draw[line width=0.4pt] (0,0) circle (0.55ex);}
% Validation numbers for the MethodsX manuscript. Every number in the text is a macro defined here.
%
% STATUS: PROVISIONAL. Values below come from the earlier assessment protocol
% (all hours with a value, lag fitted against clear-sky POA, no outage-day removal):
%   Lite, Core, PVGIS : per-site table of the earlier assessment (13 sites, Nevada 1367 excluded)
%   Plus              : outputs/thesis_vs_current/metrics.csv (same protocol, no bootstrap)
% Before submission, rerun all four series with the protocol of Section "Assessment design"
% (solarmodel.bench scoring + paired circular block bootstrap + PVGIS series) and regenerate this file.
% Plus has no bootstrap intervals yet; the counts marked (*) are point comparisons, not interval-based.

\newcommand{\nSites}{13}
\newcommand{\nSitesAll}{14}
\newcommand{\nBoot}{2000}
\newcommand{\blockDays}{5}

% medians across sites: bias (%), absolute bias (%), NRMSE (%), KGE, sites within +-5 %
\newcommand{\liteBias}{11.8}   \newcommand{\liteAbsBias}{11.8}  \newcommand{\liteNrmse}{71} \newcommand{\liteKge}{0.80} \newcommand{\liteWithinFive}{4}
\newcommand{\coreBias}{0.8}    \newcommand{\coreAbsBias}{12.4}  \newcommand{\coreNrmse}{74} \newcommand{\coreKge}{0.80} \newcommand{\coreWithinFive}{5}
\newcommand{\plusBias}{$-$1.0} \newcommand{\plusAbsBias}{9.9}   \newcommand{\plusNrmse}{73} \newcommand{\plusKge}{0.79} \newcommand{\plusWithinFive}{4}
\newcommand{\pvgisBias}{12.2}  \newcommand{\pvgisAbsBias}{14.4} \newcommand{\pvgisNrmse}{75} \newcommand{\pvgisKge}{0.77} \newcommand{\pvgisWithinFive}{4}

% benchmark: tier against PVGIS (median reduction in |bias| and NRMSE, percentage points; site counts)
\newcommand{\liteGainBias}{0.1}  \newcommand{\liteGainNrmse}{5.9} \newcommand{\liteCloserNrmse}{7}
\newcommand{\coreGainBias}{1.3}  \newcommand{\coreGainNrmse}{6.2} \newcommand{\coreCloserNrmse}{7}
\newcommand{\plusGainBias}{0.2}  \newcommand{\plusGainNrmse}{7.2} \newcommand{\plusCloserNrmse}{8}  % (*)
\newcommand{\coreCloserBias}{4}  \newcommand{\pvgisCloserBiasCore}{2}  % interval-based

% Plus against Core (*)
\newcommand{\plusCoreDeltaBias}{$-$1.9}
\newcommand{\plusCoreDeltaNrmse}{$-$0.3}
\newcommand{\plusCoreSmallerAbsBias}{6}
\newcommand{\plusCoreLowerNrmse}{9}

% Las Vegas single-site studies (solarmodel.bench protocol, studies/LasVegas_34)
\newcommand{\lvFixHours}{3{,}949}
\newcommand{\lvCoreBias}{+0.2}       \newcommand{\lvCoreNrmse}{21.0}
\newcommand{\lvPlusBias}{$-$1.2}     \newcommand{\lvPlusNrmse}{20.9}
\newcommand{\lvIamDelta}{$-$1.1}     \newcommand{\lvPressureDelta}{$-$0.3}
\newcommand{\lvSrcHours}{10{,}010}
\newcommand{\lvSrcCoreEraBias}{+4.2}  \newcommand{\lvSrcCoreEraNrmse}{18.0}
\newcommand{\lvSrcCoreNsrdbBias}{+0.6}\newcommand{\lvSrcCoreNsrdbNrmse}{14.6}
\newcommand{\lvSrcCoreMidcBias}{+1.9} \newcommand{\lvSrcCoreMidcNrmse}{14.5}
\newcommand{\lvSrcPlusEraBias}{+2.9}  \newcommand{\lvSrcPlusEraNrmse}{17.7}
\newcommand{\lvSrcPlusNsrdbBias}{$-$1.0}\newcommand{\lvSrcPlusNsrdbNrmse}{14.8}
\newcommand{\lvSrcPlusMidcBias}{+0.3} \newcommand{\lvSrcPlusMidcNrmse}{14.4}
\newcommand{\lvSandiaMidcBias}{+0.6}  \newcommand{\lvSandiaMidcNrmse}{14.4}
\newcommand{\lvInHours}{11{,}303}
\newcommand{\lvMosNrmseCut}{1.4}

% verification
\newcommand{\fitResidual}{$10^{-8}$}

% ---- Evidence classes (provisional; final classification depends on the provenance of U_c, U_v)
% A: metadata geometry and installed module          -> Las Vegas
% B: metadata geometry, surrogate module             -> Arbuckle, Golden, Henderson, Social Circle (assessment restricted to PVDAQ; Lausanne dropped)
% C: geometry inferred from measured output          -> Clayton 25, Clayton 27, Braga
% D: tracker or unrecorded/assumed tilt modelled as fixed -> Boulder, Cary SF1, Cary SF2, Calgary
\newcommand{\nClassA}{1} \newcommand{\nClassB}{4} \newcommand{\nClassC}{3} \newcommand{\nClassD}{4}
\newcommand{\nPrimary}{5}
% medians over the primary set (classes A and B, 6 sites): bias, |bias|, NRMSE, KGE
\newcommand{\liteBiasP}{13.1}  \newcommand{\liteAbsBiasP}{13.1}  \newcommand{\liteNrmseP}{51} \newcommand{\liteKgeP}{0.84}
\newcommand{\coreBiasP}{1.9}   \newcommand{\coreAbsBiasP}{1.9}   \newcommand{\coreNrmseP}{48} \newcommand{\coreKgeP}{0.93}
\newcommand{\plusBiasP}{0.3}   \newcommand{\plusAbsBiasP}{1.4}   \newcommand{\plusNrmseP}{48} \newcommand{\plusKgeP}{0.93}
\newcommand{\pvgisBiasP}{13.6} \newcommand{\pvgisAbsBiasP}{13.6} \newcommand{\pvgisNrmseP}{59} \newcommand{\pvgisKgeP}{0.82}
        % verification constants (\fitResidual)
% generated by make_validation_tables.py; do not edit
% Every assessment number quoted in main.tex and supplement.tex.
% Sites: 29 (Albuquerque_1429, Arbuckle_2107, CherryHill_1201, CherryHill_1202, Chester_1204, Cocoa_1403, Gaithersburg_4902, Gaithersburg_4903, Golden_10, Golden_1208, Golden_1283, Golden_1289, Golden_1432, Golden_1433, Golden_33, Golden_4, Golden_50, Golden_51, Henderson_1423, HuntValley_1199, LasVegas_1276, LasVegas_1277, LasVegas_1418, LasVegas_1420, LasVegas_34, Linthicum_1200, PresqueIsle_1239, Williston_1422, Wilmington_1203).
\newcommand{\vNScreened}{158}   % PVDAQ systems screened
\newcommand{\vNCandidates}{35}   % candidates after the metadata screen
\newcommand{\vNExcludedLate}{6}   % candidates not configured or excluded after extraction
\newcommand{\vNConfigured}{29}   % assessment sites (configured, not excluded)
\newcommand{\vNSites}{29}   % assessment sites with outputs
\newcommand{\vNInd}{25}   % independent validation sites
\newcommand{\vNCal}{4}   % calibrated sites
\newcommand{\vNClassA}{20}
\newcommand{\vNClassB}{5}
\newcommand{\vNStation}{16}   % independent sites with station irradiance
\newcommand{\vNStationAll}{16}
\newcommand{\vNStationGolden}{10}
\newcommand{\vNStationLV}{6}
\newcommand{\vNHours}{107{,}601}   % scored hours, all sites
\newcommand{\vHoursMin}{3{,}378}
\newcommand{\vHoursMax}{3{,}980}
\newcommand{\vDcMin}{1.0}
\newcommand{\vDcMax}{893}
\newcommand{\vYearMin}{2011}
\newcommand{\vYearMax}{2022}
\newcommand{\vAltMax}{1{,}995}
\newcommand{\vDegMin}{0.895}
\newcommand{\vDegMax}{0.998}
\newcommand{\vNSdm}{2}
\newcommand{\vReps}{2{,}000}
\newcommand{\vBlock}{5}
\newcommand{\vSeed}{20261109}
\newcommand{\vMedRatEraBias}{+15.3}
\newcommand{\vMedRatEraAbsBias}{15.3}
\newcommand{\vMedRatEraNrmse}{38.6}
\newcommand{\vMedRatEraKge}{0.79}
\newcommand{\vFiveRatEra}{3}
\newcommand{\vTenRatEra}{6}
\newcommand{\vMedRatNsrdbBias}{+12.3}
\newcommand{\vMedRatNsrdbAbsBias}{12.3}
\newcommand{\vMedRatNsrdbNrmse}{29.2}
\newcommand{\vMedRatNsrdbKge}{0.84}
\newcommand{\vFiveRatNsrdb}{5}
\newcommand{\vTenRatNsrdb}{11}
\newcommand{\vMedModEraBias}{+9.7}
\newcommand{\vMedModEraAbsBias}{9.7}
\newcommand{\vMedModEraNrmse}{37.5}
\newcommand{\vMedModEraKge}{0.83}
\newcommand{\vFiveModEra}{8}
\newcommand{\vTenModEra}{13}
\newcommand{\vMedModNsrdbBias}{+8.1}
\newcommand{\vMedModNsrdbAbsBias}{8.1}
\newcommand{\vMedModNsrdbNrmse}{27.5}
\newcommand{\vMedModNsrdbKge}{0.86}
\newcommand{\vFiveModNsrdb}{9}
\newcommand{\vTenModNsrdb}{14}
\newcommand{\vMedSysEraBias}{+6.7}
\newcommand{\vMedSysEraAbsBias}{6.7}
\newcommand{\vMedSysEraNrmse}{37.4}
\newcommand{\vMedSysEraKge}{0.84}
\newcommand{\vFiveSysEra}{12}
\newcommand{\vTenSysEra}{16}
\newcommand{\vMedSysNsrdbBias}{+2.6}
\newcommand{\vMedSysNsrdbAbsBias}{6.0}
\newcommand{\vMedSysNsrdbNrmse}{27.5}
\newcommand{\vMedSysNsrdbKge}{0.87}
\newcommand{\vFiveSysNsrdb}{10}
\newcommand{\vTenSysNsrdb}{18}
\newcommand{\vMedPvgisBias}{+6.5}
\newcommand{\vMedPvgisAbsBias}{6.5}
\newcommand{\vMedPvgisNrmse}{38.2}
\newcommand{\vMedPvgisKge}{0.82}
\newcommand{\vFivePvgis}{11}
\newcommand{\vTenPvgis}{15}
\newcommand{\vNPvgis}{25}
\newcommand{\vCalSysEraBias}{+4.6}
\newcommand{\vCalSysEraAbsBias}{4.6}
\newcommand{\vCalSysEraNrmse}{35.1}
\newcommand{\vCalSysEraKge}{0.85}
\newcommand{\vCalSysNsrdbBias}{+7.0}
\newcommand{\vCalSysNsrdbAbsBias}{7.0}
\newcommand{\vCalSysNsrdbNrmse}{24.9}
\newcommand{\vCalSysNsrdbKge}{0.90}
\newcommand{\vCalPvgisBias}{+2.6}
\newcommand{\vCalPvgisAbsBias}{4.3}
\newcommand{\vCalPvgisNrmse}{36.2}
\newcommand{\vCalPvgisKge}{0.85}
\newcommand{\vCalRatNsrdbBias}{+13.8}
\newcommand{\vCalRatNsrdbAbsBias}{13.8}
\newcommand{\vCalRatNsrdbNrmse}{28.1}
\newcommand{\vCalRatNsrdbKge}{0.84}
\newcommand{\vCalModNsrdbBias}{+9.5}
\newcommand{\vCalModNsrdbAbsBias}{9.5}
\newcommand{\vCalModNsrdbNrmse}{25.7}
\newcommand{\vCalModNsrdbKge}{0.88}
\newcommand{\vStnSysEraBias}{+11.0}
\newcommand{\vStnSysEraAbsBias}{11.0}
\newcommand{\vStnSysEraNrmse}{39.5}
\newcommand{\vStnSysEraKge}{0.79}
\newcommand{\vStnSysNsrdbBias}{+0.9}
\newcommand{\vStnSysNsrdbAbsBias}{5.3}
\newcommand{\vStnSysNsrdbNrmse}{30.0}
\newcommand{\vStnSysNsrdbKge}{0.86}
\newcommand{\vStnSysStnBias}{$-$0.5}
\newcommand{\vStnSysStnAbsBias}{3.8}
\newcommand{\vStnSysStnNrmse}{22.0}
\newcommand{\vStnSysStnKge}{0.89}
\newcommand{\vStnPvgisBias}{+12.6}
\newcommand{\vStnPvgisAbsBias}{12.6}
\newcommand{\vStnPvgisNrmse}{40.3}
\newcommand{\vStnPvgisKge}{0.75}
\newcommand{\vStnRatNsrdbBias}{+10.1}
\newcommand{\vStnRatNsrdbAbsBias}{10.1}
\newcommand{\vStnRatNsrdbNrmse}{31.3}
\newcommand{\vStnRatNsrdbKge}{0.84}
\newcommand{\vStnModNsrdbBias}{+4.4}
\newcommand{\vStnModNsrdbAbsBias}{5.4}
\newcommand{\vStnModNsrdbNrmse}{30.3}
\newcommand{\vStnModNsrdbKge}{0.86}
\newcommand{\vPairRatModEraAbsBetter}{20}
\newcommand{\vPairRatModEraAbsWorse}{2}
\newcommand{\vPairRatModEraAbsTie}{3}
\newcommand{\vPairRatModEraAbsMed}{+4.8}
\newcommand{\vPairRatModEraNrmseBetter}{14}
\newcommand{\vPairRatModEraNrmseWorse}{3}
\newcommand{\vPairRatModEraNrmseTie}{8}
\newcommand{\vPairRatModEraNrmseMed}{+1.0}
\newcommand{\vPairRatModEraKgeBetter}{22}
\newcommand{\vPairRatModEraKgeWorse}{1}
\newcommand{\vPairRatModEraKgeTie}{2}
\newcommand{\vPairRatModEraKgeMed}{+0.03}
\newcommand{\vPairRatModEraN}{25}
\newcommand{\vPairModSysEraAbsBetter}{18}
\newcommand{\vPairModSysEraAbsWorse}{1}
\newcommand{\vPairModSysEraAbsTie}{6}
\newcommand{\vPairModSysEraAbsMed}{+2.9}
\newcommand{\vPairModSysEraNrmseBetter}{11}
\newcommand{\vPairModSysEraNrmseWorse}{2}
\newcommand{\vPairModSysEraNrmseTie}{12}
\newcommand{\vPairModSysEraNrmseMed}{+0.9}
\newcommand{\vPairModSysEraKgeBetter}{15}
\newcommand{\vPairModSysEraKgeWorse}{4}
\newcommand{\vPairModSysEraKgeTie}{6}
\newcommand{\vPairModSysEraKgeMed}{+0.02}
\newcommand{\vPairModSysEraN}{25}
\newcommand{\vPairRatModNsrdbAbsBetter}{17}
\newcommand{\vPairRatModNsrdbAbsWorse}{2}
\newcommand{\vPairRatModNsrdbAbsTie}{6}
\newcommand{\vPairRatModNsrdbAbsMed}{+2.9}
\newcommand{\vPairRatModNsrdbNrmseBetter}{11}
\newcommand{\vPairRatModNsrdbNrmseWorse}{4}
\newcommand{\vPairRatModNsrdbNrmseTie}{10}
\newcommand{\vPairRatModNsrdbNrmseMed}{+0.7}
\newcommand{\vPairRatModNsrdbKgeBetter}{18}
\newcommand{\vPairRatModNsrdbKgeWorse}{1}
\newcommand{\vPairRatModNsrdbKgeTie}{6}
\newcommand{\vPairRatModNsrdbKgeMed}{+0.02}
\newcommand{\vPairRatModNsrdbN}{25}
\newcommand{\vPairModSysNsrdbAbsBetter}{14}
\newcommand{\vPairModSysNsrdbAbsWorse}{4}
\newcommand{\vPairModSysNsrdbAbsTie}{7}
\newcommand{\vPairModSysNsrdbAbsMed}{+1.9}
\newcommand{\vPairModSysNsrdbNrmseBetter}{7}
\newcommand{\vPairModSysNsrdbNrmseWorse}{4}
\newcommand{\vPairModSysNsrdbNrmseTie}{14}
\newcommand{\vPairModSysNsrdbNrmseMed}{+0.2}
\newcommand{\vPairModSysNsrdbKgeBetter}{13}
\newcommand{\vPairModSysNsrdbKgeWorse}{4}
\newcommand{\vPairModSysNsrdbKgeTie}{8}
\newcommand{\vPairModSysNsrdbKgeMed}{+0.01}
\newcommand{\vPairModSysNsrdbN}{25}
\newcommand{\vPairRatModStnAbsBetter}{11}
\newcommand{\vPairRatModStnAbsWorse}{0}
\newcommand{\vPairRatModStnAbsTie}{5}
\newcommand{\vPairRatModStnAbsMed}{+2.8}
\newcommand{\vPairRatModStnNrmseBetter}{8}
\newcommand{\vPairRatModStnNrmseWorse}{3}
\newcommand{\vPairRatModStnNrmseTie}{5}
\newcommand{\vPairRatModStnNrmseMed}{+1.1}
\newcommand{\vPairRatModStnKgeBetter}{9}
\newcommand{\vPairRatModStnKgeWorse}{1}
\newcommand{\vPairRatModStnKgeTie}{6}
\newcommand{\vPairRatModStnKgeMed}{+0.01}
\newcommand{\vPairRatModStnN}{16}
\newcommand{\vPairModSysStnAbsBetter}{8}
\newcommand{\vPairModSysStnAbsWorse}{4}
\newcommand{\vPairModSysStnAbsTie}{4}
\newcommand{\vPairModSysStnAbsMed}{+1.9}
\newcommand{\vPairModSysStnNrmseBetter}{4}
\newcommand{\vPairModSysStnNrmseWorse}{3}
\newcommand{\vPairModSysStnNrmseTie}{9}
\newcommand{\vPairModSysStnNrmseMed}{+0.1}
\newcommand{\vPairModSysStnKgeBetter}{7}
\newcommand{\vPairModSysStnKgeWorse}{4}
\newcommand{\vPairModSysStnKgeTie}{5}
\newcommand{\vPairModSysStnKgeMed}{+0.01}
\newcommand{\vPairModSysStnN}{16}
\newcommand{\vPairEraNsrdbAbsBetter}{10}
\newcommand{\vPairEraNsrdbAbsWorse}{7}
\newcommand{\vPairEraNsrdbAbsTie}{8}
\newcommand{\vPairEraNsrdbAbsMed}{+0.8}
\newcommand{\vPairEraNsrdbNrmseBetter}{25}
\newcommand{\vPairEraNsrdbNrmseWorse}{0}
\newcommand{\vPairEraNsrdbNrmseTie}{0}
\newcommand{\vPairEraNsrdbNrmseMed}{+7.8}
\newcommand{\vPairEraNsrdbKgeBetter}{21}
\newcommand{\vPairEraNsrdbKgeWorse}{0}
\newcommand{\vPairEraNsrdbKgeTie}{4}
\newcommand{\vPairEraNsrdbKgeMed}{+0.05}
\newcommand{\vPairEraNsrdbN}{25}
\newcommand{\vPairNsrdbStnAbsBetter}{5}
\newcommand{\vPairNsrdbStnAbsWorse}{3}
\newcommand{\vPairNsrdbStnAbsTie}{8}
\newcommand{\vPairNsrdbStnAbsMed}{+0.4}
\newcommand{\vPairNsrdbStnNrmseBetter}{10}
\newcommand{\vPairNsrdbStnNrmseWorse}{2}
\newcommand{\vPairNsrdbStnNrmseTie}{4}
\newcommand{\vPairNsrdbStnNrmseMed}{+2.5}
\newcommand{\vPairNsrdbStnKgeBetter}{11}
\newcommand{\vPairNsrdbStnKgeWorse}{1}
\newcommand{\vPairNsrdbStnKgeTie}{4}
\newcommand{\vPairNsrdbStnKgeMed}{+0.01}
\newcommand{\vPairNsrdbStnN}{16}
\newcommand{\vPairEraPvgisAbsBetter}{5}
\newcommand{\vPairEraPvgisAbsWorse}{9}
\newcommand{\vPairEraPvgisAbsTie}{11}
\newcommand{\vPairEraPvgisAbsMed}{$-$0.5}
\newcommand{\vPairEraPvgisNrmseBetter}{8}
\newcommand{\vPairEraPvgisNrmseWorse}{6}
\newcommand{\vPairEraPvgisNrmseTie}{11}
\newcommand{\vPairEraPvgisNrmseMed}{$-$0.3}
\newcommand{\vPairEraPvgisKgeBetter}{9}
\newcommand{\vPairEraPvgisKgeWorse}{9}
\newcommand{\vPairEraPvgisKgeTie}{7}
\newcommand{\vPairEraPvgisKgeMed}{0.00}
\newcommand{\vPairEraPvgisN}{25}
\newcommand{\vPairPvgisNsrdbAbsBetter}{10}
\newcommand{\vPairPvgisNsrdbAbsWorse}{6}
\newcommand{\vPairPvgisNsrdbAbsTie}{9}
\newcommand{\vPairPvgisNsrdbAbsMed}{+0.3}
\newcommand{\vPairPvgisNsrdbNrmseBetter}{21}
\newcommand{\vPairPvgisNsrdbNrmseWorse}{0}
\newcommand{\vPairPvgisNsrdbNrmseTie}{4}
\newcommand{\vPairPvgisNsrdbNrmseMed}{+7.7}
\newcommand{\vPairPvgisNsrdbKgeBetter}{17}
\newcommand{\vPairPvgisNsrdbKgeWorse}{4}
\newcommand{\vPairPvgisNsrdbKgeTie}{4}
\newcommand{\vPairPvgisNsrdbKgeMed}{+0.06}
\newcommand{\vPairPvgisNsrdbN}{25}
\newcommand{\vStepZeroEraBias}{$-$4.83}
\newcommand{\vStepZeroEraNrmse}{$-$1.02}
\newcommand{\vStepZeroEraMin}{$-$14.2}
\newcommand{\vStepZeroEraMax}{+1.0}
\newcommand{\vStepZeroEraAct}{25}
\newcommand{\vStepZeroNsrdbBias}{$-$4.50}
\newcommand{\vStepZeroNsrdbNrmse}{$-$0.67}
\newcommand{\vStepZeroNsrdbMin}{$-$14.1}
\newcommand{\vStepZeroNsrdbMax}{+1.6}
\newcommand{\vStepZeroNsrdbAct}{25}
\newcommand{\vStepOneEraBias}{$-$0.32}
\newcommand{\vStepOneEraNrmse}{$-$0.02}
\newcommand{\vStepOneEraMin}{$-$1.0}
\newcommand{\vStepOneEraMax}{0.0}
\newcommand{\vStepOneEraAct}{18}
\newcommand{\vStepOneNsrdbBias}{$-$0.31}
\newcommand{\vStepOneNsrdbNrmse}{$-$0.01}
\newcommand{\vStepOneNsrdbMin}{$-$0.9}
\newcommand{\vStepOneNsrdbMax}{0.0}
\newcommand{\vStepOneNsrdbAct}{18}
\newcommand{\vStepTwoEraBias}{$-$1.55}
\newcommand{\vStepTwoEraNrmse}{$-$0.54}
\newcommand{\vStepTwoEraMin}{$-$2.1}
\newcommand{\vStepTwoEraMax}{$-$1.0}
\newcommand{\vStepTwoEraAct}{25}
\newcommand{\vStepTwoNsrdbBias}{$-$1.70}
\newcommand{\vStepTwoNsrdbNrmse}{$-$0.26}
\newcommand{\vStepTwoNsrdbMin}{$-$2.5}
\newcommand{\vStepTwoNsrdbMax}{$-$1.0}
\newcommand{\vStepTwoNsrdbAct}{25}
\newcommand{\vStepThreeEraBias}{0.00}
\newcommand{\vStepThreeEraNrmse}{0.00}
\newcommand{\vStepThreeEraMin}{0.0}
\newcommand{\vStepThreeEraMax}{0.0}
\newcommand{\vStepThreeEraAct}{0}
\newcommand{\vStepThreeNsrdbBias}{0.00}
\newcommand{\vStepThreeNsrdbNrmse}{0.00}
\newcommand{\vStepThreeNsrdbMin}{0.0}
\newcommand{\vStepThreeNsrdbMax}{0.0}
\newcommand{\vStepThreeNsrdbAct}{0}
\newcommand{\vStepFourEraBias}{$-$0.01}
\newcommand{\vStepFourEraNrmse}{0.00}
\newcommand{\vStepFourEraMin}{$-$0.3}
\newcommand{\vStepFourEraMax}{+0.4}
\newcommand{\vStepFourEraAct}{6}
\newcommand{\vStepFourNsrdbBias}{$-$0.01}
\newcommand{\vStepFourNsrdbNrmse}{0.00}
\newcommand{\vStepFourNsrdbMin}{$-$0.4}
\newcommand{\vStepFourNsrdbMax}{+0.3}
\newcommand{\vStepFourNsrdbAct}{4}
\newcommand{\vStepFiveEraBias}{0.00}
\newcommand{\vStepFiveEraNrmse}{0.00}
\newcommand{\vStepFiveEraMin}{$-$0.2}
\newcommand{\vStepFiveEraMax}{+0.2}
\newcommand{\vStepFiveEraAct}{4}
\newcommand{\vStepFiveNsrdbBias}{0.00}
\newcommand{\vStepFiveNsrdbNrmse}{0.00}
\newcommand{\vStepFiveNsrdbMin}{$-$0.2}
\newcommand{\vStepFiveNsrdbMax}{+0.3}
\newcommand{\vStepFiveNsrdbAct}{3}
\newcommand{\vStepSixEraBias}{0.00}
\newcommand{\vStepSixEraNrmse}{0.00}
\newcommand{\vStepSixEraMin}{$-$0.2}
\newcommand{\vStepSixEraMax}{0.0}
\newcommand{\vStepSixEraAct}{1}
\newcommand{\vStepSixNsrdbBias}{0.00}
\newcommand{\vStepSixNsrdbNrmse}{0.00}
\newcommand{\vStepSixNsrdbMin}{$-$0.1}
\newcommand{\vStepSixNsrdbMax}{0.0}
\newcommand{\vStepSixNsrdbAct}{1}
\newcommand{\vStepSevenEraBias}{0.00}
\newcommand{\vStepSevenEraNrmse}{0.00}
\newcommand{\vStepSevenEraMin}{0.0}
\newcommand{\vStepSevenEraMax}{0.0}
\newcommand{\vStepSevenEraAct}{0}
\newcommand{\vStepSevenNsrdbBias}{0.00}
\newcommand{\vStepSevenNsrdbNrmse}{0.00}
\newcommand{\vStepSevenNsrdbMin}{0.0}
\newcommand{\vStepSevenNsrdbMax}{0.0}
\newcommand{\vStepSevenNsrdbAct}{0}
\newcommand{\vStepEightEraBias}{$-$1.75}
\newcommand{\vStepEightEraNrmse}{$-$0.12}
\newcommand{\vStepEightEraMin}{$-$7.6}
\newcommand{\vStepEightEraMax}{0.0}
\newcommand{\vStepEightEraAct}{17}
\newcommand{\vStepEightNsrdbBias}{$-$1.39}
\newcommand{\vStepEightNsrdbNrmse}{0.00}
\newcommand{\vStepEightNsrdbMin}{$-$6.0}
\newcommand{\vStepEightNsrdbMax}{0.0}
\newcommand{\vStepEightNsrdbAct}{17}
\newcommand{\vStepNineEraBias}{+0.03}
\newcommand{\vStepNineEraNrmse}{+0.01}
\newcommand{\vStepNineEraMin}{0.0}
\newcommand{\vStepNineEraMax}{+1.5}
\newcommand{\vStepNineEraAct}{8}
\newcommand{\vStepNineNsrdbBias}{+0.02}
\newcommand{\vStepNineNsrdbNrmse}{0.00}
\newcommand{\vStepNineNsrdbMin}{0.0}
\newcommand{\vStepNineNsrdbMax}{+1.4}
\newcommand{\vStepNineNsrdbAct}{8}
\newcommand{\vStepTenEraBias}{0.00}
\newcommand{\vStepTenEraNrmse}{0.00}
\newcommand{\vStepTenEraMin}{0.0}
\newcommand{\vStepTenEraMax}{0.0}
\newcommand{\vStepTenEraAct}{0}
\newcommand{\vStepTenNsrdbBias}{0.00}
\newcommand{\vStepTenNsrdbNrmse}{0.00}
\newcommand{\vStepTenNsrdbMin}{0.0}
\newcommand{\vStepTenNsrdbMax}{0.0}
\newcommand{\vStepTenNsrdbAct}{0}
\newcommand{\vStepSdmEraBias}{$-$0.37}
\newcommand{\vStepSdmEraNrmse}{$-$0.38}
\newcommand{\vStepSdmEraMin}{$-$1.0}
\newcommand{\vStepSdmEraMax}{+0.2}
\newcommand{\vStepSdmEraAct}{2}
\newcommand{\vStepSdmNsrdbBias}{$-$0.50}
\newcommand{\vStepSdmNsrdbNrmse}{$-$0.57}
\newcommand{\vStepSdmNsrdbMin}{$-$1.2}
\newcommand{\vStepSdmNsrdbMax}{+0.2}
\newcommand{\vStepSdmNsrdbAct}{2}
\newcommand{\vSysNsrdbBiasLo}{$-$11.7}
\newcommand{\vSysNsrdbBiasHi}{+13.6}
\newcommand{\vPathRatNrmseGainEra}{2.3}
\newcommand{\vPathModNrmseGainEra}{0.9}
\newcommand{\vInputNrmseGain}{7.8}
\newcommand{\vInputNrmseGainLo}{1.2}
\newcommand{\vInputNrmseGainHi}{21.4}
\newcommand{\vNInputBeatsPath}{22}
\newcommand{\vNGolden}{10}
\newcommand{\vGoldenEraBiasLo}{$-$1.6}
\newcommand{\vGoldenEraBiasHi}{+27.2}
\newcommand{\vGoldenEraBiasMed}{+20.1}
\newcommand{\vGoldenNsrdbBiasLo}{$-$11.7}
\newcommand{\vGoldenNsrdbBiasHi}{+12.1}
\newcommand{\vGoldenNsrdbBiasMed}{+5.8}
\newcommand{\vGoldenStnBiasLo}{$-$11.1}
\newcommand{\vGoldenStnBiasHi}{+13.7}
\newcommand{\vGoldenStnBiasMed}{+5.0}
\newcommand{\vGoldenPvgisBiasLo}{$-$2.5}
\newcommand{\vGoldenPvgisBiasHi}{+36.3}
\newcommand{\vNDesert}{6}
\newcommand{\vDesertEraBiasLo}{$-$4.1}
\newcommand{\vDesertEraBiasHi}{+3.7}
\newcommand{\vDesertEraBiasMed}{$-$0.3}
\newcommand{\vDesertNsrdbBiasLo}{$-$6.7}
\newcommand{\vDesertNsrdbBiasHi}{+2.6}
\newcommand{\vDesertNsrdbBiasMed}{$-$1.6}
\newcommand{\vDesertStnBiasLo}{$-$6.3}
\newcommand{\vDesertStnBiasHi}{+1.1}
\newcommand{\vDesertStnBiasMed}{$-$2.3}
\newcommand{\vDesertPvgisBiasLo}{$-$1.1}
\newcommand{\vDesertPvgisBiasHi}{+4.5}
\newcommand{\vNEast}{7}
\newcommand{\vEastEraBiasLo}{$-$3.9}
\newcommand{\vEastEraBiasHi}{+16.7}
\newcommand{\vEastEraBiasMed}{+3.4}
\newcommand{\vEastNsrdbBiasLo}{$-$3.1}
\newcommand{\vEastNsrdbBiasHi}{+13.6}
\newcommand{\vEastNsrdbBiasMed}{+8.4}
\newcommand{\vEastPvgisBiasLo}{$-$3.6}
\newcommand{\vEastPvgisBiasHi}{+13.3}
\newcommand{\vHorizonMax}{9.9}
\newcommand{\vNHorizon}{29}
\newcommand{\vPvgisBeatsSysEraNrmse}{11}
\newcommand{\vNsrdbBeatsPvgisNrmse}{25}
\newcommand{\vNsrdbBeatsEraNrmse}{25}
\newcommand{\vAblPresModEra}{$-$0.32}
\newcommand{\vAblPresModEraLo}{$-$1.0}
\newcommand{\vAblPresModEraHi}{0.0}
\newcommand{\vAblPresModEraAct}{18}
\newcommand{\vAblPresSysEra}{$-$0.32}
\newcommand{\vAblPresSysEraLo}{$-$1.1}
\newcommand{\vAblPresSysEraHi}{0.0}
\newcommand{\vAblPresSysEraAct}{18}
\newcommand{\vAblPresModNsrdb}{$-$0.31}
\newcommand{\vAblPresModNsrdbLo}{$-$0.9}
\newcommand{\vAblPresModNsrdbHi}{0.0}
\newcommand{\vAblPresModNsrdbAct}{18}
\newcommand{\vAblPresSysNsrdb}{$-$0.31}
\newcommand{\vAblPresSysNsrdbLo}{$-$1.0}
\newcommand{\vAblPresSysNsrdbHi}{0.0}
\newcommand{\vAblPresSysNsrdbAct}{18}
\newcommand{\vAblIamModEra}{$-$1.55}
\newcommand{\vAblIamModEraLo}{$-$2.1}
\newcommand{\vAblIamModEraHi}{$-$1.0}
\newcommand{\vAblIamModEraAct}{25}
\newcommand{\vAblIamSysEra}{$-$1.69}
\newcommand{\vAblIamSysEraLo}{$-$2.1}
\newcommand{\vAblIamSysEraHi}{$-$1.0}
\newcommand{\vAblIamSysEraAct}{25}
\newcommand{\vAblIamModNsrdb}{$-$1.70}
\newcommand{\vAblIamModNsrdbLo}{$-$2.5}
\newcommand{\vAblIamModNsrdbHi}{$-$1.0}
\newcommand{\vAblIamModNsrdbAct}{25}
\newcommand{\vAblIamSysNsrdb}{$-$1.86}
\newcommand{\vAblIamSysNsrdbLo}{$-$2.4}
\newcommand{\vAblIamSysNsrdbHi}{$-$1.1}
\newcommand{\vAblIamSysNsrdbAct}{25}
\newcommand{\vAblHorModEra}{0.00}
\newcommand{\vAblHorModEraLo}{0.0}
\newcommand{\vAblHorModEraHi}{0.0}
\newcommand{\vAblHorModEraAct}{0}
\newcommand{\vAblHorSysEra}{0.00}
\newcommand{\vAblHorSysEraLo}{0.0}
\newcommand{\vAblHorSysEraHi}{0.0}
\newcommand{\vAblHorSysEraAct}{0}
\newcommand{\vAblHorModNsrdb}{0.00}
\newcommand{\vAblHorModNsrdbLo}{0.0}
\newcommand{\vAblHorModNsrdbHi}{0.0}
\newcommand{\vAblHorModNsrdbAct}{0}
\newcommand{\vAblHorSysNsrdb}{0.00}
\newcommand{\vAblHorSysNsrdbLo}{0.0}
\newcommand{\vAblHorSysNsrdbHi}{0.0}
\newcommand{\vAblHorSysNsrdbAct}{0}
\newcommand{\vAblLossModEra}{$-$0.01}
\newcommand{\vAblLossModEraLo}{$-$0.4}
\newcommand{\vAblLossModEraHi}{+0.6}
\newcommand{\vAblLossModEraAct}{9}
\newcommand{\vAblLossSysEra}{$-$0.02}
\newcommand{\vAblLossSysEraLo}{$-$0.1}
\newcommand{\vAblLossSysEraHi}{+0.4}
\newcommand{\vAblLossSysEraAct}{6}
\newcommand{\vAblLossModNsrdb}{$-$0.01}
\newcommand{\vAblLossModNsrdbLo}{$-$0.5}
\newcommand{\vAblLossModNsrdbHi}{+0.4}
\newcommand{\vAblLossModNsrdbAct}{6}
\newcommand{\vAblLossSysNsrdb}{$-$0.02}
\newcommand{\vAblLossSysNsrdbLo}{$-$0.1}
\newcommand{\vAblLossSysNsrdbHi}{+0.3}
\newcommand{\vAblLossSysNsrdbAct}{4}
\newcommand{\vAblSoilModEra}{+0.01}
\newcommand{\vAblSoilModEraLo}{$-$0.2}
\newcommand{\vAblSoilModEraHi}{+0.2}
\newcommand{\vAblSoilModEraAct}{5}
\newcommand{\vAblSoilSysEra}{0.00}
\newcommand{\vAblSoilSysEraLo}{$-$0.2}
\newcommand{\vAblSoilSysEraHi}{+0.3}
\newcommand{\vAblSoilSysEraAct}{5}
\newcommand{\vAblSoilModNsrdb}{0.00}
\newcommand{\vAblSoilModNsrdbLo}{$-$0.2}
\newcommand{\vAblSoilModNsrdbHi}{+0.2}
\newcommand{\vAblSoilModNsrdbAct}{4}
\newcommand{\vAblSoilSysNsrdb}{0.00}
\newcommand{\vAblSoilSysNsrdbLo}{$-$0.2}
\newcommand{\vAblSoilSysNsrdbHi}{+0.3}
\newcommand{\vAblSoilSysNsrdbAct}{4}
\newcommand{\vAblClipModEra}{0.00}
\newcommand{\vAblClipModEraLo}{$-$0.5}
\newcommand{\vAblClipModEraHi}{0.0}
\newcommand{\vAblClipModEraAct}{2}
\newcommand{\vAblClipSysEra}{0.00}
\newcommand{\vAblClipSysEraLo}{$-$0.2}
\newcommand{\vAblClipSysEraHi}{0.0}
\newcommand{\vAblClipSysEraAct}{1}
\newcommand{\vAblClipModNsrdb}{0.00}
\newcommand{\vAblClipModNsrdbLo}{$-$0.7}
\newcommand{\vAblClipModNsrdbHi}{0.0}
\newcommand{\vAblClipModNsrdbAct}{3}
\newcommand{\vAblClipSysNsrdb}{0.00}
\newcommand{\vAblClipSysNsrdbLo}{$-$0.1}
\newcommand{\vAblClipSysNsrdbHi}{0.0}
\newcommand{\vAblClipSysNsrdbAct}{1}
\newcommand{\vAblDepthModEra}{0.00}
\newcommand{\vAblDepthModEraLo}{0.0}
\newcommand{\vAblDepthModEraHi}{0.0}
\newcommand{\vAblDepthModEraAct}{0}
\newcommand{\vAblDepthSysEra}{0.00}
\newcommand{\vAblDepthSysEraLo}{0.0}
\newcommand{\vAblDepthSysEraHi}{0.0}
\newcommand{\vAblDepthSysEraAct}{0}
\newcommand{\vAblDepthModNsrdb}{0.00}
\newcommand{\vAblDepthModNsrdbLo}{0.0}
\newcommand{\vAblDepthModNsrdbHi}{0.0}
\newcommand{\vAblDepthModNsrdbAct}{0}
\newcommand{\vAblDepthSysNsrdb}{0.00}
\newcommand{\vAblDepthSysNsrdbLo}{0.0}
\newcommand{\vAblDepthSysNsrdbHi}{0.0}
\newcommand{\vAblDepthSysNsrdbAct}{0}
\newcommand{\vAblTrigModEra}{$-$1.76}
\newcommand{\vAblTrigModEraLo}{$-$7.8}
\newcommand{\vAblTrigModEraHi}{0.0}
\newcommand{\vAblTrigModEraAct}{17}
\newcommand{\vAblTrigSysEra}{$-$1.75}
\newcommand{\vAblTrigSysEraLo}{$-$7.6}
\newcommand{\vAblTrigSysEraHi}{0.0}
\newcommand{\vAblTrigSysEraAct}{17}
\newcommand{\vAblTrigModNsrdb}{$-$1.41}
\newcommand{\vAblTrigModNsrdbLo}{$-$6.1}
\newcommand{\vAblTrigModNsrdbHi}{0.0}
\newcommand{\vAblTrigModNsrdbAct}{17}
\newcommand{\vAblTrigSysNsrdb}{$-$1.39}
\newcommand{\vAblTrigSysNsrdbLo}{$-$6.0}
\newcommand{\vAblTrigSysNsrdbHi}{0.0}
\newcommand{\vAblTrigSysNsrdbAct}{17}
\newcommand{\vAblAlbModEra}{+0.10}
\newcommand{\vAblAlbModEraLo}{0.0}
\newcommand{\vAblAlbModEraHi}{+1.6}
\newcommand{\vAblAlbModEraAct}{12}
\newcommand{\vAblAlbSysEra}{+0.03}
\newcommand{\vAblAlbSysEraLo}{0.0}
\newcommand{\vAblAlbSysEraHi}{+1.5}
\newcommand{\vAblAlbSysEraAct}{8}
\newcommand{\vAblAlbModNsrdb}{+0.07}
\newcommand{\vAblAlbModNsrdbLo}{0.0}
\newcommand{\vAblAlbModNsrdbHi}{+1.5}
\newcommand{\vAblAlbModNsrdbAct}{10}
\newcommand{\vAblAlbSysNsrdb}{+0.02}
\newcommand{\vAblAlbSysNsrdbLo}{0.0}
\newcommand{\vAblAlbSysNsrdbHi}{+1.4}
\newcommand{\vAblAlbSysNsrdbAct}{8}
\newcommand{\vAblSnowModEra}{$-$0.11}
\newcommand{\vAblSnowModEraLo}{$-$10.0}
\newcommand{\vAblSnowModEraHi}{0.0}
\newcommand{\vAblSnowModEraAct}{13}
\newcommand{\vAblSnowSysEra}{$-$2.08}
\newcommand{\vAblSnowSysEraLo}{$-$17.5}
\newcommand{\vAblSnowSysEraHi}{0.0}
\newcommand{\vAblSnowSysEraAct}{17}
\newcommand{\vAblSnowModNsrdb}{$-$0.10}
\newcommand{\vAblSnowModNsrdbLo}{$-$10.2}
\newcommand{\vAblSnowModNsrdbHi}{0.0}
\newcommand{\vAblSnowModNsrdbAct}{13}
\newcommand{\vAblSnowSysNsrdb}{$-$2.11}
\newcommand{\vAblSnowSysNsrdbLo}{$-$15.8}
\newcommand{\vAblSnowSysNsrdbHi}{0.0}
\newcommand{\vAblSnowSysNsrdbAct}{17}
\newcommand{\vBmsDistMin}{0.3}
\newcommand{\vBmsDistMax}{1}
\newcommand{\vUnlvDistMin}{6.6}
\newcommand{\vUnlvDistMax}{22}
\newcommand{\vUcFit}{1}
\newcommand{\vUcNoct}{3}
\newcommand{\vUcClass}{25}
    % assessment numbers, written by make_validation_tables.py

\externaldocument{build/main}
\renewcommand{\thetable}{S\arabic{table}}
\renewcommand{\thesection}{S\arabic{section}}
\renewcommand{\theHtable}{S\arabic{table}}
\renewcommand{\theHsection}{S\arabic{section}}

\begin{document}

\begin{center}
{\Large Supplementary material}\\[4pt]
\RePV: An open, configuration-driven physical modelling pipeline for photovoltaic power estimation
\end{center}

Equation, table and section numbers without the prefix S, and all citations, refer to the main article.

%===========================================================================
\section{Method specification}
\label{sec:s_spec}

Table~\ref{tab:s_options} defines \RePV{} and the simple rating model of the validation stage by stage. Tables~\ref{tab:weather} and~\ref{tab:config} give the weather field mapping and the configuration groups, Table~\ref{tab:s_prov} the provenance classes and Table~\ref{tab:outputs} the output series of one run. The key-by-key schema with accepted values and ranges is \texttt{solarmodel/schema.py} of the repository; in the code, \RePV{} is the model \texttt{repv} and the simple rating model is \texttt{rating}.

\begin{table}[htbp]
\centering
\caption{Stages of \RePV{} and of the simple rating model used as a reference in the validation. Both run on the same weather and the same plane-of-array geometry.}
\label{tab:s_options}
\small
\begin{tabularx}{\textwidth}{>{\raggedright\arraybackslash}p{0.2\textwidth} >{\raggedright\arraybackslash}X >{\raggedright\arraybackslash}p{0.27\textwidth}}
\toprule
Stage & \RePV & Simple rating model \\
\midrule
System data & module datasheet or CEC entry, U-values, static losses; optional row layout, ILR, inverter entry & DC and AC rating, $\gamma_P$, NOCT \\
Solar geometry & SPA at the interval midpoint, site-pressure air mass & as \RePV{} \\
Transposition & Reindl, isotropic ground reflection & as \RePV{} \\
Terrain horizon & PVGIS horizon; beam and circumsolar blocked & none \\
Row geometry & infinite sheds, applied as ratios, when a layout is given & none \\
Incidence angle & Fresnel, beam and diffuse & none \\
Spectrum & SAPM air-mass polynomial & none \\
Snow & NREL coverage with 6~h trigger and ERA5-Land depth condition and reset, applied to the light; snow albedo & none \\
Soiling & Kimber, on the effective irradiance & none \\
Cell temperature & PVsyst heat balance (Eq.~\ref{eq:tcell}) & NOCT model with efficiency correction~\cite{duffieBeckman1991} \\
DC power & De Soto single diode & linear in $G_{\mathrm{POA}}$ with $\gamma_P$~\cite{dobosPVWattsVersion52014} \\
Losses & mismatch, DC wiring and age before the inverter, AC wiring after it & one combined factor after conversion \\
Inverter & PVWatts at $P_{\mathrm{ac0}}$; Sandia optional & included in the combined factor \\
AC limit & site AC capacity & site AC capacity \\
\bottomrule
\end{tabularx}
\end{table}

\begin{table}[htbp]
\centering
\caption{Weather inputs as retrieved from the Open-Meteo archive API. Temporal meanings follow the API documentation~\cite{zippenfenig2023openmeteo}.}
\label{tab:weather}
\small
\begin{tabularx}{\textwidth}{>{\raggedright\arraybackslash}p{0.09\textwidth} >{\raggedright\arraybackslash}p{0.30\textwidth} >{\raggedright\arraybackslash}p{0.13\textwidth} >{\raggedright\arraybackslash}X >{\raggedright\arraybackslash}p{0.13\textwidth}}
\toprule
Symbol & API field & Product & Temporal meaning & Conversion \\
\midrule
$G_{\mathrm{GHI}}$ & \texttt{shortwave\_radiation} & ERA5 & mean of the preceding hour & none \\
$G_{\mathrm{DNI}}$ & \texttt{direct\_normal\_irradiance} & ERA5 & mean of the preceding hour & none \\
$G_{\mathrm{DHI}}$ & \texttt{diffuse\_radiation} & ERA5 & mean of the preceding hour & none \\
$T_a$ & \texttt{temperature\_2m} & API default & instantaneous at the label & none \\
$v_{10}$ & \texttt{wind\_speed\_10m} & API default & instantaneous at the label & requested in m\,s$^{-1}$ \\
$S$ & \texttt{snowfall} & ERA5 & sum over the preceding hour, as fresh-snow depth & none (cm) \\
$Z$ & \texttt{snow\_depth} & ERA5-Land & instantaneous at the label & m to cm \\
$P_{\mathrm{d}}$ & \texttt{precipitation\_sum} (daily) & ERA5 & sum over the UTC day & none (mm) \\
\bottomrule
\end{tabularx}
\end{table}

\begin{table}[htbp]
\centering
\caption{Configuration groups of a site file. A missing required entry, an unknown section or key, a value outside its range and an unknown provenance class stop the run with an error that names the entry.}
\label{tab:config}
\small
\begin{tabularx}{\textwidth}{>{\raggedright\arraybackslash}p{0.19\textwidth} >{\raggedright\arraybackslash}X >{\raggedright\arraybackslash}p{0.2\textwidth}}
\toprule
Group & Principal entries & Status \\
\midrule
Location and array & latitude, longitude, altitude, tilt, azimuth, albedo; tracking and tracker geometry; snow tilt & required; tracker if tracking \\
System & modules per string, strings, DC rating, site AC capacity & modules and strings required \\
Module & STC values, temperature coefficients, cell count; CEC entry, band gap, bifaciality & STC values, coefficients and cell count required \\
Optics & refractive index $n$, extinction $K$, glass thickness $L$ & required \\
Thermal & $U_c$, $U_v$, $\eta_{\mathrm{th}}$, $\alpha_{\mathrm{abs}}$; NOCT block for the rating model & required \\
Losses & mismatch, DC and AC wiring, inverter efficiency; age factor & required; age factor optional \\
Soiling & loss rate, cleaning threshold, grace period, cap, wash dates & optional; Kimber defaults \\
Layout and inverter & ground coverage ratio, height, collector width; ILR; CEC inverter entry and count & optional \\
Snow & removal: sliding or fast & optional; sliding \\
Provenance & one class per parameter (Table~\ref{tab:s_prov}) & complete for every assessment site \\
\bottomrule
\end{tabularx}
\end{table}

\begin{table}[htbp]
\centering
\caption{Provenance classes of the site files and their use at the \vNConfigured{} assessment sites.}
\label{tab:s_prov}
\small
\begin{tabularx}{\textwidth}{l >{\raggedright\arraybackslash}X}
\toprule
Class & Use at the assessment sites \\
\midrule
\texttt{measured} & $U_c$ fitted to on-site module temperature sensors (Henderson 1423) \\
\texttt{datasheet} & module STC values and temperature coefficients from the CEC entry or datasheet; $\eta_{\mathrm{th}}$; $U_c$ from the datasheet NOCT \\
\texttt{published} & optical constants of the glass, PVsyst $U_c$ classes, $U_v=0$, $\alpha_{\mathrm{abs}}$, Kimber soiling values, age degradation rate \\
\texttt{metadata} & coordinates, altitude, orientation, DC rating, module and inverter counts, site AC capacity \\
\texttt{output} & fast snow removal at the \vNCal{} calibrated sites \\
\texttt{assumed} & albedo, static loss factors, string split without string data, rating-model factors \\
\bottomrule
\end{tabularx}
\end{table}

\begin{table}[htbp]
\centering
\caption{Output series of one run, on the weather index.}
\label{tab:outputs}
\small
\begin{tabularx}{\textwidth}{l >{\raggedright\arraybackslash}X l}
\toprule
Column & Quantity & Unit \\
\midrule
\texttt{p\_ac\_kw} & AC power at the grid connection point & kW \\
\texttt{p\_dc\_kw} & array maximum-power DC output before mismatch, DC-wiring and age losses & kW \\
\texttt{clip\_kw} & inverter clipping plus site AC capacity limitation, AC side & kW \\
\texttt{poa\_global}, \texttt{poa\_eff} & $G_{\mathrm{POA}}$ before snow; $R_{\mathrm{soil}}G_{\mathrm{eff,s}}$ & W\,m$^{-2}$ \\
\texttt{t\_cell} & cell temperature & $^{\circ}$C \\
\texttt{snow\_tx}, \texttt{soiling} & $R_{\mathrm{snow}}$ and $R_{\mathrm{soil}}$ & $-$ \\
\bottomrule
\end{tabularx}
\end{table}

The output \texttt{clip\_kw} reports the power removed by the inverter ceiling, referred to the AC side through $\eta_{\mathrm{ac}}$, plus the power removed by the site AC capacity. For instantaneous irradiance inputs such as PVGIS-SARAH, solar position is evaluated at the stated instant instead of the interval midpoint.

\subsection{Module parameter extraction}
\label{sec:s_fit}

\paragraph{Datasheet fit.} Reference parameters are fitted to the datasheet values $I_{\mathrm{sc}}$, $V_{\mathrm{oc}}$, $I_{\mathrm{mp}}$, $V_{\mathrm{mp}}$, $\alpha_{\mathrm{sc}}$, $\beta_{\mathrm{oc}}$ and the cell count $N_{\mathrm{cell}}$ with \texttt{pvlib.ivtools.sdm.fit\_desoto}, with $E_{g,\mathrm{ref}}=1.121$~eV unless the site file gives another value. The fit starts from $I_{L,0}=I_{\mathrm{sc}}$, $a_0=nN_{\mathrm{cell}}kT_{\mathrm{ref}}/q$, $I_{0,0}=I_{L,0}\exp(-V_{\mathrm{oc}}/a_0)$, $R_{s,0}=0.3(V_{\mathrm{oc}}-V_{\mathrm{mp}})/I_{\mathrm{mp}}$ and $R_{\mathrm{sh},0}=m_{\mathrm{sh}}\max[V_{\mathrm{mp}}/(I_{\mathrm{sc}}-I_{\mathrm{mp}}),300~\Omega]$, and is attempted over $n\in\{1.1,1.0,1.2,1.3,1.5\}$ and $m_{\mathrm{sh}}\in\{1,2,5,10,0.5\}$ in that order. \texttt{fit\_desoto} solves the five equations implied by these datasheet conditions; the first converged set in this fixed order is used, which makes the extraction deterministic, and if no combination converges the run stops with an error. The verification suite requires $I_{\mathrm{sc}}$, $V_{\mathrm{oc}}$ and $P_{\mathrm{mp}}$ at STC to reproduce the datasheet within a relative $10^{-8}$.

\paragraph{CEC entry.} Where the site file names a CEC database entry, its reference parameters $a_{\mathrm{ref}}$, $I_{L,\mathrm{ref}}$, $I_{0,\mathrm{ref}}$, $R_s$ and $R_{\mathrm{sh,ref}}$ are used, $\alpha_{\mathrm{sc}}$ is the datasheet value, and the sixth CEC parameter (Adjust) is not used, because it trades the current and voltage coefficients against the power coefficient. $E_{g,\mathrm{ref}}$ is then the root in $[0.5,2.5]$~eV of $\mathrm{d}V_{\mathrm{oc}}/\mathrm{d}T=\beta_{\mathrm{oc}}$ at 1000~W\,m$^{-2}$ and 25~$^{\circ}$C, so the open-circuit voltage follows the datasheet coefficient whatever the cell technology; the run stops with an error if no root exists. This path is used where the datasheet point has only fits with a negative shunt resistance, and for every module with a CEC entry in the validation. At the assessment sites, the solved $E_{g,\mathrm{ref}}$ lay between 1.01 and 1.10~eV for the silicon modules and was 1.26~eV for the CIGS module of Golden 1289. The verification suite requires the CEC STC curve within a relative $10^{-6}$ and $\mathrm{d}V_{\mathrm{oc}}/\mathrm{d}T$ within a relative $10^{-4}$ of $\beta_{\mathrm{oc}}$.

\subsection{Snow state: implementation details}
\label{sec:s_snow}

\RePV{} passes the trailing 6~h snowfall sum, the ERA5-Land snow depth and the plane-of-array irradiance to \texttt{pvlib.snow.coverage\_nrel}, with its published coefficients and depth threshold of 1~cm. The removal condition of Eq.~\ref{eq:snow_reset} lets snow slide when the air temperature exceeds 0~$^{\circ}$C at zero irradiance, a threshold that falls by 1~K per 80~W\,m$^{-2}$ of plane-of-array irradiance. The state $C_t$ of Eq.~\ref{eq:snow_state} applies to the output of hour $t$, with $C=0$ before the first hour of the run. An event takes precedence over sliding in the same hour, and because an event requires $Z_t\geq1$~cm, an event and the depth reset cannot coincide. With the 6~h trigger, every hour whose trailing 6~h snowfall exceeds 1~cm is an event, so coverage is held at 1 until that sum falls to the threshold. Without ground snow depth, sliding would be the only removal term, so a snow tilt $\beta_{\mathrm{s}}=0$ would retain the first cover indefinitely. The snow tilt equals the array tilt for fixed arrays and $\beta(t)$ for trackers; the run stops with an error for a fixed array with zero snow tilt unless a separate non-zero snow tilt is supplied. The fast-removal setting replaces the cover by zero and keeps the snow albedo.

\subsection{Worked example}
\label{sec:s_example}

Configuration file of the Las Vegas system (PVDAQ 34), without the comments and the provenance table that accompany every value in the repository copy:
{\scriptsize
\begin{verbatim}
[location]  lat = 36.1952   lon = -115.1582   altitude = 620.0
[array]     tracking = false   tilt = 11.2   azimuth = 180.0   albedo = 0.20
[system]    rated_capacity_kw = 146.64   ac_cap_kw = 135.0   ns_arr = 13   np_arr = 47
[module]    isc_ref = 8.79   voc_ref = 37.00   imp_ref = 7.85   vmp_ref = 30.60
            mu_isc = 0.00387   mu_voc = -0.1217   ns = 60   eg_ref = 1.121
[iam]       n = 1.526   k = 4.0   l = 0.002
[thermal.pvsyst]  u_c = 28.3   u_v = 0.0   eta_module = 0.1465   alpha_abs = 0.9
[losses]    n_degradation = 0.9957   n_mismatch = 0.98   n_dc_wiring = 0.98
            n_ac_wiring = 0.99   n_inverter = 0.96
[inverter]  cec = "Satcon_Technology__PVS_135__480V_"   count = 1
[provenance]  "array.tilt" = "metadata"   "array.albedo" = "assumed"
              "thermal.pvsyst.u_c" = "datasheet"   ...
\end{verbatim}}

The tables below are generated by \texttt{python make\_worked\_example.py --site LasVegas\_34 --year 2011}, which calls the same functions as the assessment. The hourly excerpt lists the weather inputs, the intermediate series and the AC power for one day, so that a reimplementation can be checked stage by stage.

\begin{table}[htbp]
\centering
\caption{Worked example, hourly excerpt: PVDAQ 34, \RePV{}, local day 2011-06-21 (labels in UTC, hour ending). Inputs from Table~\ref{tab:weather}, intermediate series and output from Table~\ref{tab:outputs}. Generated by \texttt{make\_worked\_example.py}.}
\label{tab:worked_hours}
\small
\setlength{\tabcolsep}{4pt}
\begin{tabular}{l rrr r rr r rr}
\toprule
Hour & $G_{\mathrm{GHI}}$ & $G_{\mathrm{DNI}}$ & $G_{\mathrm{DHI}}$ & $T_a$ & $G_{\mathrm{POA}}$ & \texttt{poa\_eff} & $T_c$ & $P_{\mathrm{dc}}$ & $P_{\mathrm{ac}}$ \\
(UTC) & \multicolumn{3}{c}{[W\,m$^{-2}$]} & [$^{\circ}$C] & \multicolumn{2}{c}{[W\,m$^{-2}$]} & [$^{\circ}$C] & \multicolumn{2}{c}{[kW]} \\
\midrule
13:00 & 13 & 89 & 7 & 22.9 & 6 & 5 & 23.1 & 0.6 & 0.0 \\
14:00 & 146 & 510 & 47 & 24.6 & 104 & 63 & 27.4 & 8.6 & 7.1 \\
15:00 & 347 & 712 & 72 & 27.7 & 310 & 244 & 36.1 & 33.6 & 30.3 \\
16:00 & 554 & 822 & 87 & 30.1 & 533 & 451 & 44.6 & 61.0 & 55.5 \\
17:00 & 741 & 883 & 99 & 31.8 & 739 & 635 & 51.8 & 83.7 & 76.3 \\
18:00 & 891 & 916 & 110 & 33.3 & 904 & 778 & 57.8 & 100.0 & 91.1 \\
19:00 & 990 & 937 & 113 & 33.8 & 1014 & 871 & 61.3 & 110.1 & 100.3 \\
20:00 & 1034 & 945 & 116 & 35.0 & 1063 & 911 & 63.8 & 113.9 & 103.7 \\
21:00 & 1017 & 943 & 114 & 35.9 & 1044 & 895 & 64.2 & 111.7 & 101.8 \\
22:00 & 940 & 926 & 112 & 36.8 & 958 & 823 & 62.8 & 103.4 & 94.2 \\
23:00 & 810 & 898 & 105 & 37.4 & 813 & 700 & 59.5 & 89.2 & 81.3 \\
00:00 & 638 & 852 & 93 & 37.6 & 623 & 532 & 54.5 & 69.1 & 62.9 \\
01:00 & 437 & 766 & 80 & 37.4 & 404 & 331 & 48.4 & 43.7 & 39.6 \\
02:00 & 228 & 608 & 60 & 36.7 & 184 & 129 & 41.7 & 16.9 & 14.8 \\
03:00 & 53 & 287 & 22 & 35.4 & 17 & 13 & 35.9 & 1.5 & 0.6 \\
\bottomrule
\end{tabular}
\end{table}

\begin{table}[htbp]
\centering
\caption{Worked example, year 2011: simulated AC energy, specific yield and capacity factor (on 146.64~kWp DC) of \RePV{} and the simple rating model, ERA5 input.}
\label{tab:worked_year}
\small
\begin{tabular}{l rrr}
\toprule
Model & Energy [MWh] & Specific yield [kWh\,kWp$^{-1}$] & Capacity factor [\%] \\
\midrule
Simple rating model & 281.5 & 1920 & 21.9 \\
\RePV & 254.9 & 1738 & 19.8 \\
\bottomrule
\end{tabular}
\end{table}


%===========================================================================
\section{Sites}
\label{sec:s_sites}

% generated by make_validation_tables.py; do not edit
\begin{table}[htbp]
\centering
\caption{Site-selection flow. The basis column states what each decision used: metadata only, the measured series only, or the measured series compared with model output. The measured-series exclusions rest on checks that use no model output, except where stated.}
\label{tab:s_selection}
\footnotesize
\begin{tabularx}{\textwidth}{>{\raggedright\arraybackslash}X >{\raggedright\arraybackslash}p{0.22\textwidth} r}
\toprule
Step & Basis & Systems \\
\midrule
PVDAQ systems screened (OEDI system list and Data Prize folder) &  & 158 \\
\quad excluded: no AC power series in the record & metadata & $-$78 \\
\quad excluded: net-export meter only & metadata & $-$30 \\
\quad excluded: single-axis tracker & metadata & $-$2 \\
\quad excluded: tilt or azimuth not recorded & metadata & $-$11 \\
\quad excluded: two orientations on one AC series & metadata & $-$1 \\
\quad excluded: diagnostic clock offset at the $\pm4$~h search bound & measured series & $-$1 \\
Candidates: fixed single-plane array, AC series, module recorded &  & 35 \\
\quad not configured: PVDAQ 36, modules with built-in reflectors, which the chain does not represent & metadata & $-$1 \\
\quad excluded: New Smyrna Beach 1231, battery-based inverter, so the AC series follows the battery and the load, not the array & metadata & $-$1 \\
\quad excluded: New Orleans 1275, diurnal profile inconsistent with the recorded azimuth of 135$^{\circ}$ and output about 30\,\% below the recorded rating throughout the day, consistent with unrecorded shading & measured series and model output & $-$1 \\
\quad excluded: Golden 1332, the documented 1.15~MWp is split between a garage roof and a south facade in an undocumented ratio, and the chain represents one plane & metadata & $-$1 \\
\quad excluded: Social Circle 9069, the measured series sums an undocumented subset of the plant's AC meters & measured series & $-$1 \\
\quad excluded: Golden 1284, hourly peak 1.10 to 1.15 times the DC rating from the module records, with no source for the installed rating & measured series & $-$1 \\
Assessment sites &  & 29 \\
\quad of which calibrated (snow removal selected on the scored year) & measured series and model output & 4 \\
\quad of which independent validation &  & 25 \\
\bottomrule
\end{tabularx}
\end{table}


% generated by make_validation_tables.py; do not edit
\begin{table}[htbp]
\centering
\caption{Assessment sites. Tilt and azimuth (degrees from north), DC rating and the recorded module are from the PVDAQ metadata (DC rating: module count times the wattage in the module model name, the recorded system power where the module records lack a wattage, or the documented rating where one is published); class A uses the database entry of the recorded module, class B a surrogate or an assumed variant (Table~\ref{tab:s_classb}); c: calibrated (snow removal selected on the scored year). Hours: scored hours. Station: MIDC station used as the station irradiance input. $f_{\mathrm{deg}}$: age degradation factor. Clock: whole-hour shift applied by the extraction in the scored year, standard time; daylight-saving time.}
\label{tab:s_sites}
\scriptsize
\setlength{\tabcolsep}{1.5pt}
\makebox[\textwidth][c]{\begin{tabular}{l r l r >{\raggedright\arraybackslash}p{0.23\textwidth} l r r l r l}
\toprule
Site & Alt. (m) & Tilt/az. & DC (kWp) & Recorded module (technology, count) & Cl. & Year & Hours & Station & $f_{\mathrm{deg}}$ & Clock (h) \\
\midrule
Albuquerque 1429 (NM) & 1658 & 35/180 & 5.9 & suniva opt270-60 blk-blk (mono-Si, 22) & A & 2017 & 3{,}424 & n/a & 0.993 & 0; 0 \\
Arbuckle 2107 (CA) & 10 & 25/180 & 893.0 & Hyundai HiS-M310TI (mono-Si, 120) & B & 2019 & 3{,}899 & n/a & 0.988 & 0; 0 \\
Cherry Hill 1201 (NJ) & 18 & 10/230 & 140.1 & SolarWorld sw245 mono (mono-Si, 572) & A$^{\mathrm{c}}$ & 2012 & 3{,}970 & n/a & 0.993 & $-$1; $-$1 \\
Cherry Hill 1202 (NJ) & 19 & 10/230 & 51.8 & SolarWorld sw240 (mono-Si, 216) & B & 2013 & 3{,}980 & n/a & 0.988 & $-$1; $-$1 \\
Chester 1204 (PA) & 10 & 5/140 & 78.6 & SolarWorld sw245 mono (mono-Si, 321) & A$^{\mathrm{c}}$ & 2014 & 3{,}952 & n/a & 0.983 & 0; $-$1 \\
Cocoa 1403 (FL) & 12 & 35/180 & 5.9 & suniva opt270-60 blk-blk (mono-Si, 22) & A & 2016 & 3{,}890 & n/a & 0.993 & 0; 0 \\
Gaithersburg 4902 (MD) & 138 & 20/180 & 270.7 & Sharp NU-U235F2 (mono-Si, 1152) & A$^{\mathrm{c}}$ & 2017 & 3{,}856 & n/a & 0.986 & 0; 0 \\
Gaithersburg 4903 (MD) & 149 & 10/180 & 73.3 & Sharp NU-U235F2 (mono-Si, 312) & A & 2016 & 3{,}764 & n/a & 0.990 & 0; 0 \\
Golden 4 (CO) & 1795 & 40/180 & 1.0 & Sanyo hip-200ba3 (mono-Si, 5) & B & 2015 & 3{,}535 & BMS & 0.961 & 0; +1 \\
Golden 10 (CO) & 1793 & 40/180 & 1.1 & shell solar eclipse 80 (cis family thin-film, 1) & A & 2021 & 3{,}803 & BMS & 0.925 & 0; 0 \\
Golden 33 (CO) & 1794 & 40/180 & 2.4 & silicor materials 240 (Unknown, 1) & B & 2022 & 3{,}782 & BMS & 0.939 & 0; 0/+1 \\
Golden 50 (CO) & 1995 & 45/158 & 7.4 & siemens M55 c-SI (Unknown, 1) & A & 2016 & 3{,}635 & BMS & 0.895 & 0/+1; +1 \\
Golden 51 (CO) & 1995 & 45/158 & 7.4 & siemens M55 c-Si (Unknown, 1) & A & 2012 & 3{,}658 & BMS & 0.913 & 0; 0/+1 \\
Golden 1208 (CO) & 1780 & 9.1/165 & 524.2 & SunPower spr-315e-wht-d (mono-Si, 1664) & A & 2017 & 3{,}543 & BMS & 0.971 & 0; 0/+1 \\
Golden 1283 (CO) & 1800 & 10/165 & 408.2 & SunPower spr-315e-wht-d (mono-Si, 1296) & A & 2020 & 3{,}518 & BMS & 0.958 & 0; 0/+1 \\
Golden 1289 (CO) & 1794 & 40/180 & 1.8 & stion corporation stn-130 (cylindrical cigs, 14) & A & 2016 & 3{,}690 & BMS & 0.981 & 0/+1; 0 \\
Golden 1432 (CO) & 1780 & 10/175 & 94.0 & Evergreen es-190-sl (ribbon polycrystalline si, 495) & A & 2012 & 3{,}605 & BMS & 0.986 & $-$1; $-$1/0 \\
Golden 1433 (CO) & 1780 & 10/180 & 383.8 & Sanyo hit-205a (SJT, 936) & A & 2011 & 3{,}378 & BMS & 0.998 & $-$1; $-$1 \\
Henderson 1423 (NV) & 697 & 35/180 & 6.5 & suniva opt270-60 blk-blk (mono-Si, 24) & A & 2022 & 3{,}522 & UNLV & 0.967 & 0; 0 \\
Hunt Valley 1199 (MD) & 325 & 20/180 & 52.9 & Suntech 280W (multi-Si, 189) & B & 2018 & 3{,}883 & n/a & 0.960 & 0; $-$1 \\
Las Vegas 34 (NV) & 620 & 11.2/180 & 146.6 & Sharp nu-u240f1 (mono-Si, 611) & A & 2011 & 3{,}583 & UNLV & 0.996 & 0; 0 \\
Las Vegas 1276 (NV) & 620 & 5/180 & 60.5 & Sharp nu-u240f1 (mono-Si, 252) & A & 2016 & 3{,}748 & UNLV & 0.976 & 0; 0 \\
Las Vegas 1277 (NV) & 620 & 10/180 & 40.6 & Sharp nu-u240f1 (mono-Si, 169) & A & 2012 & 3{,}510 & UNLV & 0.996 & 0; 0 \\
Las Vegas 1418 (NV) & 779 & 10/138 & 43.7 & SolarWorld sw280 mono (mono-Si, 156) & A & 2016 & 3{,}467 & UNLV & 0.995 & 0; $-$1 \\
Las Vegas 1420 (NV) & 678 & 10/180 & 35.8 & SunPower spr-320e-wht-d (mono-Si, 112) & A & 2019 & 3{,}727 & UNLV & 0.980 & 0; $-$1 \\
Linthicum 1200 (MD) & 155 & 10/205 & 51.8 & SolarWorld sw240 (mono-Si, 216) & B$^{\mathrm{c}}$ & 2017 & 3{,}937 & n/a & 0.967 & 0; $-$1/0 \\
Presque Isle 1239 (ME) & 197 & 10/155 & 20.2 & Sharp nd-224uc1 (multi-Si, 90) & A & 2015 & 3{,}546 & n/a & 0.981 & 0; 0 \\
Williston 1422 (VT) & 109 & 35/180 & 6.5 & suniva opt270-60 blk-blk (mono-Si, 24) & A & 2015 & 3{,}819 & n/a & 0.997 & 0; 0 \\
Wilmington 1203 (DE) & 4 & 20/205 & 197.5 & SolarWorld sw245 mono (mono-Si, 806) & A & 2016 & 3{,}977 & n/a & 0.973 & $-$1; $-$2/$-$1 \\
\bottomrule
\end{tabular}}
\end{table}

\begin{table}[htbp]
\centering
\caption{Class-B sites: surrogate module or assumed variant.}
\label{tab:s_classb}
\footnotesize
\begin{tabularx}{\textwidth}{l X}
\toprule
Site & Module parameters used \\
\midrule
Arbuckle 2107 & STC curve of the HiS-M310KI for the recorded HiS-M310TI (De Soto fit of the TI entry has no positive shunt resistance) \\
Cherry Hill 1202 & mono variant assumed for the recorded `sw240' \\
Golden 33 & 60-cell 240~W module for the recorded Silicor 240 (no public datasheet) \\
Golden 4 & CEC HIP-200BA20 for the recorded HIP-200BA3 (no BA3 entry) \\
Hunt Valley 1199 & STP280-24/Vd assumed for the recorded `Suntech 280W' \\
Linthicum 1200 & mono variant assumed for the recorded `sw240' \\
\bottomrule
\end{tabularx}
\end{table}


The DC rating of the two Golden test arrays PVDAQ 50 and 51 is the STC rating of 7.43~kWp documented in NREL/CP-411-20430, because the recorded system power of 6.0~kW is a PTC AC estimate. Two periods were left out of the extraction. At Golden 50, the record from 2017 was left out, because the peak of one of its two half-arrays fell by about 22\,\% that year. At Golden 1208, the record from 4 August 2018 was left out, because no channel records AC power over long stretches after that date. Both scored years precede these periods.

%===========================================================================
\section{Assessment data processing}
\label{sec:s_qc}

This section gives the rules summarised in Sections~\ref{sec:design} and~\ref{sec:sites} in full. The implementation and its tests are part of the software archive.

\subsection{Site configuration}

Each site was configured from its PVDAQ metadata: coordinates, altitude, tilt, azimuth, module model and count, and inverter model and count. The DC rating was the module count multiplied by the wattage in the module model name. Where the module records give no wattage, the recorded system power was used. Module parameters were taken from the CEC or Sandia database entry of the recorded module. The heat-loss coefficient $U_c$ was fitted to on-site module temperature sensors at \vUcFit{} site, derived from the datasheet NOCT at \vUcNoct{} sites, and taken from the PVsyst value of the mounting~\cite{pvsystUvalueMeasurement2026} at \vUcClass{} sites, with the mounting inferred from the system name and tilt. Output was reduced by the age factor $f_{\mathrm{deg}}=(1-0.005)^{a}$ on the DC side (Eq.~\ref{eq:ac}), with $a$ the age in years at the middle of the scored year, counted from the earlier of the recorded installation date and the first measured value. The rate of 0.5\,\%/yr is the median field degradation rate reported by Jordan and Kurtz~\cite{jordanKurtz2013degradation}. The factor ranged from \vDegMin{} to \vDegMax. One year was scored per site, chosen from the measured series and the irradiance coverage alone: at most sites the calendar year with the most daylight hours that passed quality control (Section~\ref{sec:s_qc_extract}), and at four sites a year selected for complete coverage of the record. Every parameter of a site file carries its provenance class (Table~\ref{tab:s_prov}). Module parameters came from the CEC database entry of the recorded module, or of the surrogate of Table~\ref{tab:s_classb}, at 22 sites and from the datasheet fit of Section~\ref{sec:s_fit} at Golden 4, 10, 50, 51, 1432 and 1433 and Las Vegas 34. Each run started at the beginning of the scored year with no snow cover and a clean array.

The snow-cover ratios of the calibrated sites (Section~\ref{sec:sites}) refer to days retained by the quality control, which removes fully covered days with zero output (Section~\ref{sec:s_qc_extract}), so they describe days on which the arrays were at least partly clear.

\subsection{Extraction and quality control of the measured series}
\label{sec:s_qc_extract}

The measured series of each site was extracted from the raw PVDAQ record by rules that use the metadata, the measured values and clear-sky irradiance, and no output of \RePV{} or PVGIS. Every AC power column was averaged onto a 5-min grid. Each column was converted to kW separately. The record was split at lasting steps of the daily peak level, such as the change of units in many PVDAQ records on 4 August 2018. Each segment was scaled by the power of ten that placed its 99th-percentile daily peak nearest to 0.9 times the expected peak. Days on which the night level approached the daytime peak were set missing, because the column then records another quantity.

A whole-system series was formed month by month from the candidates: each system-total column, the sum of the inverter columns and the sum of the meter columns. A candidate counted as whole-system in a month when its peak level reached 0.9 times that of the highest candidate. The most complete whole-system candidate was used, and the others filled its gaps. Months from April to October whose level fell below 0.6 times the usual level of the calendar month were set missing as partial-array months. Months from November to March were kept, because snow cover produces the same signature, and the outage rule below decides on them. Samples above 1.3 times the DC rating or below $-5$\,\% of it were set missing, and smaller negative values, which are standby consumption, were set to zero. Runs of 30 or more days with a daily peak below 0.1\,\% of the DC rating were set missing as logger failures.

The 5-min values were averaged to hourly means. Clock offsets were corrected in blocks of one calendar month and one daylight-saving status, following the US daylight-saving calendar of the site's time zone. For each block, the whole-hour shift $\ell\in\{-4,\dots,4\}$~h was selected that maximised the Pearson correlation between the shifted series and Ineichen clear-sky GHI~\cite{ineichenPerez2002} over hours with positive clear-sky GHI. A block kept the median shift of its four neighbouring blocks unless its own optimum raised the correlation by at least 0.01. All labels were thereby moved onto fixed standard time, and from there onto UTC by the standard-time offset of the site's time zone. Hours with zero clear-sky GHI were set to zero below 2\,\% of the DC rating, as sensor offsets, and set missing above it. A month was set missing when, after the correction, its energy in hours with zero clear-sky GHI exceeded 1\,\% of its total, or when its own optimum shift still differed from zero by a correlation gain of at least 0.01. Every resulting series passed a separate check of its schema, power range, dark hours and clock in each year. The shifts applied in the scored year are listed in Table~\ref{tab:s_sites}.

\subsection{Irradiance inputs and PVGIS}
\label{sec:s_irr}

NSRDB and station irradiance were checked against the BSRN `extremely rare' limits and, where the diffuse component is measured, the closure test, as applied by Long and Shi~\cite{longShi2008qcrad}; failing values were set missing without gap filling. Station minutes were averaged to 10-min steps that contained at least 9 valid minutes. NSRDB samples were interpolated to 10-min steps through the clear-sky index. Both were then averaged to hour-ending means, an hour being formed only from six valid steps, and the chain was run on these hourly means exactly as on ERA5. The PVGIS values were taken as hourly means centred on their half-hour stamps.

\subsection{Scoring set}
\label{sec:s_scoring}

Measured hourly means were centred on their interval and compared with the simulated mean of the same hour. A local day was removed as an outage day when, during hours with clear-sky GHI above 200~W\,m$^{-2}$, it contained any missing value or three or more hours of zero output, or when its maximum output stayed below 2\,\% of the DC rating on a day with clear-sky GHI above 500~W\,m$^{-2}$. A fully snow-covered array produces the same signature, so such days are removed with the outages. Scored hours were daylight hours, with clear-sky GHI above 50~W\,m$^{-2}$, on retained days with a measured value. An hour was also removed when any compared series of the site had no value, which occurs where the satellite or station irradiance has a gap, so all series of a site were scored on the same hours against the same measured series.

%===========================================================================


%===========================================================================
\section{Results}
\label{sec:s_results}

Tables~\ref{tab:s_scores} and~\ref{tab:s_scores_nrmse} give the scores of every site with 95\,\% paired block-bootstrap intervals. Every series was run with ERA5 temperature and wind, whatever its irradiance input.

% generated by make_validation_tables.py; do not edit
\begin{table}[htbp]
\centering
\caption{Signed bias $b$ (\%) per site, with 95\,\% block-bootstrap intervals for \RePV{} on each irradiance input and for PVGIS-ERA5; the simple rating model on ERA5 and \RePV{} with the Sandia inverter model on NSRDB (sites with a CEC inverter entry) without intervals. Paired differences are in \texttt{outputs/pvdaq\_validation/bootstrap.parquet} of the archived release.}
\label{tab:s_scores}
\scriptsize
\setlength{\tabcolsep}{1.4pt}
\makebox[\textwidth][c]{\begin{tabular}{l r llll r}
\toprule
Site & Rating ERA5 & \RePV{} ERA5 & \RePV{} NSRDB & \RePV{} station & PVGIS-ERA5 & Sandia \\
\midrule
\multicolumn{7}{l}{\textit{Independent validation sites}} \\
Albuquerque 1429 & +1.1 & $-$1.0 [$-$3.2, +1.3] & $-$6.0 [$-$7.1, $-$4.9] & n/a & $-$5.0 [$-$6.9, $-$2.8] & n/a \\
Arbuckle 2107 & +15.3 & +9.0 [+6.5, +11.5] & +11.0 [+8.3, +13.6] & n/a & +6.5 [+4.8, +8.4] & +9.7 \\
Cherry Hill 1202 & +12.9 & +3.4 [+1.2, +5.9] & +3.1 [+1.5, +4.8] & n/a & +3.4 [+0.6, +6.7] & n/a \\
Cocoa 1403 & +3.3 & +2.2 [+0.3, +4.3] & +8.4 [+7.3, +9.7] & n/a & $-$3.6 [$-$5.4, $-$1.7] & n/a \\
Gaithersburg 4903 & +14.2 & +6.7 [+4.5, +9.0] & +10.2 [+8.7, +12.0] & n/a & +2.7 [+0.2, +5.2] & n/a \\
Golden 4 & +31.5 & +27.2 [+22.3, +32.8] & +5.1 [+2.9, +7.3] & +6.3 [+4.9, +8.1] & +22.5 [+17.5, +28.4] & n/a \\
Golden 10 & +26.9 & +22.4 [+19.0, +26.2] & +12.1 [+10.5, +13.8] & +7.8 [+6.8, +8.8] & +26.3 [+22.8, +30.5] & n/a \\
Golden 33 & +24.2 & +17.8 [+14.6, +21.3] & +6.6 [+5.0, +8.3] & +3.7 [+2.2, +5.4] & +21.1 [+16.8, +26.2] & n/a \\
Golden 50 & +32.8 & +27.1 [+24.1, +30.4] & +11.8 [+9.6, +14.0] & +13.7 [+12.3, +15.0] & +36.3 [+32.7, +40.3] & n/a \\
Golden 51 & +33.8 & +26.8 [+21.6, +32.4] & +6.8 [+3.2, +10.6] & +11.1 [+7.8, +14.7] & +34.7 [+29.4, +40.6] & n/a \\
Golden 1208 & +23.5 & +12.8 [+9.7, +16.0] & $-$2.7 [$-$4.9, $-$0.7] & $-$3.2 [$-$5.3, $-$1.5] & +11.8 [+8.6, +15.9] & n/a \\
Golden 1283 & +28.4 & +13.2 [+9.1, +18.1] & +2.3 [$-$0.9, +6.4] & $-$1.1 [$-$4.3, +2.9] & +18.3 [+13.1, +24.6] & n/a \\
Golden 1289 & +28.2 & +26.5 [+23.2, +30.3] & +11.2 [+9.2, +13.2] & +12.3 [+10.9, +14.0] & +19.2 [+15.8, +23.0] & n/a \\
Golden 1432 & +24.6 & +9.3 [+6.3, +12.4] & $-$4.0 [$-$6.9, $-$1.5] & $-$1.7 [$-$4.1, +0.6] & +13.5 [+10.2, +17.2] & n/a \\
Golden 1433 & +11.4 & $-$1.6 [$-$4.6, +1.7] & $-$11.7 [$-$14.0, $-$9.3] & $-$11.1 [$-$13.4, $-$8.9] & $-$2.5 [$-$6.5, +2.8] & n/a \\
Henderson 1423 & +4.8 & $-$1.2 [$-$3.5, +1.1] & $-$2.2 [$-$4.2, $-$0.1] & $-$3.9 [$-$5.7, $-$2.0] & +0.4 [$-$0.7, +1.6] & n/a \\
Hunt Valley 1199 & +8.6 & +2.1 [$-$0.6, +5.2] & $-$1.3 [$-$3.6, +0.8] & n/a & +1.6 [$-$1.4, +5.0] & n/a \\
Las Vegas 34 & +12.6 & +2.7 [0.0, +5.5] & $-$0.5 [$-$2.5, +1.6] & +0.1 [$-$2.1, +2.4] & +4.5 [+3.1, +6.0] & $-$0.2 \\
Las Vegas 1276 & +7.8 & $-$4.1 [$-$5.6, $-$2.6] & $-$6.7 [$-$8.0, $-$5.5] & $-$6.3 [$-$7.5, $-$5.0] & $-$0.6 [$-$1.5, +0.4] & n/a \\
Las Vegas 1277 & +7.3 & $-$1.2 [$-$2.9, +0.8] & $-$5.6 [$-$6.9, $-$4.2] & $-$3.4 [$-$4.8, $-$1.9] & $-$1.1 [$-$2.2, +0.3] & n/a \\
Las Vegas 1418 & +13.5 & +3.7 [+1.7, +5.9] & +2.6 [+0.8, +4.4] & +1.1 [$-$0.6, +3.0] & +2.6 [+1.3, +4.0] & n/a \\
Las Vegas 1420 & +12.3 & +0.7 [$-$0.9, +2.4] & $-$1.0 [$-$2.3, +0.3] & $-$1.3 [$-$2.7, +0.2] & +2.6 [+1.3, +4.0] & n/a \\
Presque Isle 1239 & +20.2 & $-$3.9 [$-$9.4, +1.5] & $-$3.1 [$-$9.1, +2.5] & n/a & +8.0 [+2.3, +14.8] & n/a \\
Williston 1422 & +22.6 & +16.7 [+12.6, +22.0] & +13.6 [+10.9, +16.6] & n/a & +13.3 [+7.7, +20.9] & n/a \\
Wilmington 1203 & +15.9 & +8.1 [+4.3, +12.1] & +9.2 [+6.8, +12.0] & n/a & +8.5 [+4.4, +13.0] & n/a \\
\multicolumn{7}{l}{\textit{Calibrated sites}} \\
Cherry Hill 1201 & +10.1 & +3.2 [+1.1, +5.5] & +4.1 [+3.1, +5.1] & n/a & +0.4 [$-$1.6, +2.8] & n/a \\
Chester 1204 & +18.6 & +10.9 [+7.1, +15.6] & +9.2 [+6.1, +13.3] & n/a & +7.5 [+3.9, +12.2] & n/a \\
Gaithersburg 4902 & +5.0 & +2.7 [+0.5, +5.2] & +4.7 [+3.6, +5.9] & n/a & $-$3.9 [$-$6.2, $-$1.5] & n/a \\
Linthicum 1200 & +12.4 & +6.1 [+3.4, +8.8] & +9.8 [+7.1, +12.5] & n/a & +4.8 [+2.1, +7.6] & n/a \\
\bottomrule
\end{tabular}}
\end{table}

\begin{table}[htbp]
\centering
\caption{NRMSE (\%) per site, with 95\,\% block-bootstrap intervals for \RePV{} on each irradiance input and for PVGIS-ERA5; the simple rating model on ERA5 and \RePV{} with the Sandia inverter model on NSRDB (sites with a CEC inverter entry) without intervals. Paired differences are in \texttt{outputs/pvdaq\_validation/bootstrap.parquet} of the archived release.}
\label{tab:s_scores_nrmse}
\scriptsize
\setlength{\tabcolsep}{1.4pt}
\makebox[\textwidth][c]{\begin{tabular}{l r llll r}
\toprule
Site & Rating ERA5 & \RePV{} ERA5 & \RePV{} NSRDB & \RePV{} station & PVGIS-ERA5 & Sandia \\
\midrule
\multicolumn{7}{l}{\textit{Independent validation sites}} \\
Albuquerque 1429 & 24.3 & 25.2 [21.6, 28.9] & 17.2 [14.8, 20.1] & n/a & 24.8 [22.0, 27.7] & n/a \\
Arbuckle 2107 & 28.5 & 26.8 [24.1, 29.6] & 24.4 [22.1, 26.7] & n/a & 25.6 [22.9, 28.5] & 23.3 \\
Cherry Hill 1202 & 38.1 & 35.1 [31.8, 38.9] & 25.0 [22.0, 28.4] & n/a & 36.7 [32.6, 41.7] & n/a \\
Cocoa 1403 & 31.4 & 31.3 [28.9, 33.9] & 24.2 [22.3, 26.2] & n/a & 31.6 [29.4, 33.9] & n/a \\
Gaithersburg 4903 & 33.2 & 29.8 [26.8, 32.9] & 21.3 [19.1, 23.9] & n/a & 31.2 [28.1, 34.5] & n/a \\
Golden 4 & 54.3 & 51.7 [44.7, 59.4] & 30.3 [26.0, 35.0] & 19.4 [13.3, 25.7] & 50.7 [43.2, 59.2] & n/a \\
Golden 10 & 45.9 & 44.7 [39.8, 50.1] & 27.5 [24.4, 30.8] & 17.5 [14.6, 20.4] & 46.9 [41.6, 52.7] & n/a \\
Golden 33 & 47.9 & 42.9 [38.2, 47.7] & 29.8 [26.2, 33.4] & 25.1 [20.6, 29.3] & 47.6 [41.1, 54.4] & n/a \\
Golden 50 & 61.2 & 57.3 [52.8, 61.9] & 49.2 [46.0, 52.3] & 48.4 [44.5, 51.8] & 66.7 [61.3, 72.6] & n/a \\
Golden 51 & 62.5 & 60.6 [54.6, 67.1] & 48.5 [44.6, 52.7] & 44.9 [40.4, 49.7] & 65.8 [59.1, 73.2] & n/a \\
Golden 1208 & 45.8 & 40.4 [36.4, 44.7] & 35.2 [31.1, 39.1] & 30.4 [25.4, 34.8] & 41.2 [36.4, 46.9] & n/a \\
Golden 1283 & 51.6 & 41.8 [37.1, 47.3] & 34.1 [29.3, 39.3] & 34.8 [29.7, 40.0] & 47.2 [40.8, 54.7] & n/a \\
Golden 1289 & 47.3 & 47.3 [42.9, 51.9] & 32.9 [30.2, 35.9] & 24.9 [21.8, 28.3] & 43.1 [38.7, 48.0] & n/a \\
Golden 1432 & 44.3 & 38.5 [34.4, 42.7] & 33.0 [28.8, 37.3] & 28.5 [24.0, 32.9] & 39.3 [34.9, 44.4] & n/a \\
Golden 1433 & 38.6 & 34.7 [30.7, 39.4] & 30.3 [26.6, 34.6] & 24.7 [20.0, 29.5] & 38.0 [33.6, 43.6] & n/a \\
Henderson 1423 & 17.6 & 18.9 [16.7, 21.4] & 13.8 [12.3, 15.5] & 17.1 [15.3, 19.1] & 16.5 [14.4, 18.9] & n/a \\
Hunt Valley 1199 & 38.1 & 37.4 [33.6, 41.5] & 25.5 [19.9, 31.9] & n/a & 38.2 [34.3, 43.0] & n/a \\
Las Vegas 34 & 20.5 & 19.2 [17.5, 21.0] & 15.0 [14.0, 16.2] & 15.6 [14.3, 17.1] & 17.1 [15.2, 19.1] & 14.9 \\
Las Vegas 1276 & 18.0 & 18.1 [16.5, 19.7] & 16.1 [14.9, 17.4] & 15.0 [14.0, 16.2] & 16.5 [14.6, 18.5] & n/a \\
Las Vegas 1277 & 18.0 & 17.0 [15.1, 19.1] & 15.7 [14.3, 17.3] & 14.4 [13.2, 15.5] & 16.0 [13.9, 18.3] & n/a \\
Las Vegas 1418 & 22.5 & 20.3 [18.2, 22.5] & 16.5 [15.1, 18.0] & 17.4 [16.0, 18.9] & 18.2 [16.1, 20.4] & n/a \\
Las Vegas 1420 & 19.9 & 16.7 [14.9, 18.5] & 12.7 [11.2, 14.5] & 12.8 [11.4, 14.4] & 16.1 [14.1, 18.1] & n/a \\
Presque Isle 1239 & 54.3 & 49.4 [43.2, 57.0] & 43.5 [35.3, 52.3] & n/a & 49.2 [41.8, 58.5] & n/a \\
Williston 1422 & 54.3 & 50.9 [44.3, 58.9] & 34.4 [28.8, 41.2] & n/a & 51.2 [41.6, 63.4] & n/a \\
Wilmington 1203 & 41.0 & 39.2 [34.6, 44.3] & 29.8 [25.2, 34.9] & n/a & 39.3 [34.6, 44.4] & n/a \\
\multicolumn{7}{l}{\textit{Calibrated sites}} \\
Cherry Hill 1201 & 33.3 & 31.7 [29.2, 34.5] & 20.7 [18.9, 22.6] & n/a & 32.2 [29.7, 35.1] & n/a \\
Chester 1204 & 42.2 & 39.0 [32.9, 47.2] & 29.1 [22.0, 37.8] & n/a & 39.1 [33.3, 46.9] & n/a \\
Gaithersburg 4902 & 32.0 & 31.8 [28.8, 35.1] & 20.6 [18.5, 22.9] & n/a & 33.3 [30.6, 36.2] & n/a \\
Linthicum 1200 & 40.3 & 38.4 [35.1, 42.2] & 31.3 [27.5, 35.7] & n/a & 39.5 [36.1, 43.1] & n/a \\
\bottomrule
\end{tabular}}
\end{table}


%===========================================================================
\section{Computational performance}
\label{sec:s_bench}

\begin{table}[htbp]
\centering
\caption{Wall time and memory of \RePV{} on a two-core Intel Xeon virtual machine at 2.3~GHz with 8~GB of memory (Linux, Python~3.13, pvlib~0.16.1, environment of the repository), from cached weather. Weather download time depends on the network and is not included. Median of five runs.}
\label{tab:s_bench}
\small
\begin{tabularx}{\textwidth}{>{\raggedright\arraybackslash}X r}
\toprule
Task & Value \\
\midrule
\RePV{}, one site-year at 1~h (Las Vegas 34, 2011) & 0.10~s \\
Simple rating model, one site-year & 0.09~s \\
\RePV{}, 25-year series at 1~h (219{,}144 steps) & 2.1~s \\
Complete assessment of one site: 31 to 32 model variants on two or three irradiance inputs, PVGIS and scoring & 8 to 13~s \\
Peak resident memory of the Python process & 0.4~GB \\
Share of the single-diode solution in the run time & 45\,\% \\
Share of the solar position in the run time & 25\,\% \\
\bottomrule
\end{tabularx}
\end{table}

\end{document}
